% Generated draft: requires mathematical and source review.
\documentclass[11pt]{article}
\usepackage[margin=25mm]{geometry}
\usepackage{fontspec}
\setmainfont{Latin Modern Roman}
\usepackage{amsmath,amssymb,mathtools,unicode-math}
\setmathfont{Latin Modern Math}
\usepackage{enumitem,needspace,xurl,hyperref}
\setlength{\parindent}{0pt}
\setlength{\parskip}{5pt}
\setlength{\emergencystretch}{3em}
\newcommand{\N}{\mathbb{N}}
\newcommand{\Z}{\mathbb{Z}}
\newcommand{\Q}{\mathbb{Q}}
\newcommand{\R}{\mathbb{R}}
\newcommand{\C}{\mathbb{C}}
\newcommand{\F}{\mathbb{F}}
\newcommand{\A}{\mathbb{A}}
\newcommand{\PP}{\mathbb{P}}
\newcommand{\E}{\mathbb{E}}
\newcommand{\eps}{\varepsilon}
\newcommand{\dd}{\,\mathrm{d}}
\begin{document}
\section*{UnsolvedMath ID: 30002384}
\textbf{Source problem code:} catalogue entry
\subsection{Definitions and mathematical statement}
Fix a prime \(\ell\), a number field \(K\), and an abelian variety \(A/K\). Let \(L\) range over quadratic extensions of \(K\), up to \(K\)-isomorphism; thus \([L:K]=2\). Write \(\mathcal O_L\) for the ring of integers of \(L\), and let \(\operatorname{Pic}(\mathcal O_L)\) be its ideal class group. Its \(\ell\)-torsion subgroup is
\[
\operatorname{Pic}(\mathcal O_L)[\ell]
=\{c\in\operatorname{Pic}(\mathcal O_L):\ell c=0\}.
\]
Write \(A_L=A\times_K L\) for the base change of \(A\), and \(A_L[\ell]\) for the kernel of multiplication by \(\ell\). For each place \(v\) of \(L\), let \(L_v\) be the completion and let
\[
\delta_v:A_L(L_v)/\ell A_L(L_v)\longrightarrow H^1(L_v,A_L[\ell])
\]
be the Kummer map. Here \(H^1\) denotes continuous Galois cohomology. Define the \(\ell\)-Selmer group by
\[
\operatorname{Sel}_{\ell}(A_L)
=\{\xi\in H^1(L,A_L[\ell]):
\operatorname{res}_v(\xi)\in\operatorname{im}(\delta_v)
\text{ for every place }v\text{ of }L\}.
\]
Both groups whose cardinalities occur below are finite for each fixed \(L\).

The supplied record proposes the following two unboundedness assertions, with \(\ell\), \(K\), and, in the second assertion, \(A\) fixed while \(L\) varies:
\begin{enumerate}
\item For every integer \(B\geq 1\), there exists a quadratic extension \(L/K\) such that
\[
\#\operatorname{Pic}(\mathcal O_L)[\ell]>B.
\]
\item For every integer \(B\geq 1\), there exists a quadratic extension \(L/K\) such that
\[
\#\operatorname{Sel}_{\ell}(A_L)>B.
\]
\end{enumerate}
These are separate assertions: the extensions witnessing the two inequalities need not be the same. The source presents these assertions for arbitrary prime \(\ell\), number field \(K\), and abelian variety \(A/K\). It does not explicitly require \(\dim A>0\); whether its convention excludes the zero-dimensional abelian variety is an unresolved qualification of the intended second assertion. No such hypothesis has been silently added here.
\subsection{Short English statement}
For a fixed prime and number field, the first conjecture asks whether the prime-torsion subgroup of the ideal class group can be arbitrarily large over quadratic extensions. For a fixed abelian variety over that field, the second asks the same question for its Selmer groups after base change. These are separate unboundedness conjectures. The intended dimension restriction on the abelian variety requires clarification.
\subsection{Sources}
\begin{itemize}
\item Local UnsolvedMath registry, \texttt{data/unsolvedmath/problems.json}, permanent ID \(30002384\), source problem code \texttt{catalogue entry}, title: Unbounded Torsion and Selmer Groups in Quadratic Extensions. The supplied record belongs to the Oberwolfach Reports Open Problems collection and lists proposed year 2013. No individual proposer is supplied.
\item Supplied report citation: Explicit Methods in Number Theory (2014), Oberwolfach Reports, DOI 10.4171/owr/2013/35, pp.\ 2035--2082. \url{https://doi.org/10.4171/owr/2013/35}. The record describes Oberwolfach Reports as documenting mathematical workshops and their open-problem discussions. Its citation uses 2014, while its proposed year and report identifier use 2013; this dating distinction has not been independently checked.
\item Supplied literature reference: Selmer groups and class groups. \url{https://doi.org/10.1112/S0010437X14007441}.
\item Supplied literature reference: Selmer groups and class groups. \url{https://arxiv.org/abs/1307.4261}.
\item Supplied literature reference: Oberwolfach Report 35/2013. \url{https://doi.org/10.4171/owr/2013/35}.
\item Literature assessment dated exactly \texttt{2026-08-21}: the supplied curation labels the problem partially solved and states that the class-group conjecture is known for \(\ell=2\), that the corresponding Selmer assertion is known in several \(\ell=2\) cases, and that both number-field conjectures remain open for every odd \(\ell\); it also reports additional positive cases for function-field analogues. This is supplied dated curation, not an independently verified literature assessment. The record itself requests independent verification of time-sensitive status and cited claims. No precise hypotheses for those Selmer cases or function-field analogues are supplied.
\end{itemize}
\end{document}
